\documentclass[a4paper,12pt]{article}
\usepackage[utf8]{inputenc}
\usepackage[T1]{fontenc}
\usepackage[english]{babel}
\usepackage{hyperref}
\usepackage{listings}
\usepackage{xcolor}
\usepackage{geometry}
\usepackage{graphicx}
\usepackage{booktabs}

\geometry{margin=1in}

\definecolor{codegray}{rgb}{0.5,0.5,0.5}
\definecolor{codepurple}{rgb}{0.58,0,0.82}
\definecolor{backcolour}{rgb}{0.95,0.95,0.92}

\lstdefinestyle{mystyle}{
    backgroundcolor=\color{backcolour},
    commentstyle=\color{codegray},
    keywordstyle=\color{magenta},
    numberstyle=\tiny\color{codegray},
    stringstyle=\color{codepurple},
    basicstyle=\ttfamily\footnotesize,
    breakatwhitespace=false,
    breaklines=true,
    captionpos=b,
    keepspaces=true,
    numbers=left,
    numbersep=5pt,
    showspaces=false,
    showstringspaces=false,
    showtabs=false,
    tabsize=2
}

\lstset{style=mystyle}

\title{\textbf{PASP: Portable Android Statistics Program}\\ \Large User Manual -- Version 0.3.5}
\author{Jules (Engineering AI)}
\date{\today}

\begin{document}

\maketitle
\tableofcontents
\newpage

\section{Introduction}
PASP (Portable Android Statistics Program) is a lightweight, local statistical analysis tool designed for the terminal. Inspired by the philosophy of JASP, it aims to provide a clean and intuitive interface for statistical analysis, specifically optimized for mobile environments like \textbf{Android with Termux}, but fully compatible with any modern terminal.

\section{Installation}
PASP can be installed directly from GitHub.

\subsection{Standard Installation}
\begin{lstlisting}[language=bash]
pip install git+https://github.com/ppmena/pasp
\end{lstlisting}

\subsection{Recommended for Termux (Android)}
To ensure all scientific dependencies (\texttt{numpy}, \texttt{scipy}, \texttt{pandas}) compile correctly in Termux, install the build tools first:
\begin{lstlisting}[language=bash]
pkg update && pkg upgrade
pkg install python clang make cmake pkg-config ninja meson
pip install git+https://github.com/ppmena/pasp
\end{lstlisting}

\subsection{Updating PASP}
Stay on the latest version using the built-in update command:
\begin{lstlisting}[language=bash]
pasp update
# OR
pasp --update
\end{lstlisting}

\section{Usage Styles}
PASP supports two distinct CLI styles to suit different workflows.

\subsection{Style A: Direct File Usage}
Pass the file as the first argument followed by operation flags.
\begin{lstlisting}[language=bash]
pasp data.csv --auto
pasp data.csv --ttest-ind score group --plot
\end{lstlisting}

\subsection{Style B: Subcommands}
Use specific subcommands for structured control.
\begin{lstlisting}[language=bash]
pasp descriptives data.csv
pasp anova data.csv var group
\end{lstlisting}

\section{Core Features}

\subsection{Smart Data Loading}
PASP automatically detects:
\begin{itemize}
    \item \textbf{Delimiters:} Comma (\texttt{,}), Semicolon (\texttt{;}), Tab (\texttt{\textbackslash t}), or Pipe (\texttt{|}).
    \item \textbf{Encoding:} UTF-8 or Latin-1.
\end{itemize}

\subsection{Intelligent Analysis}
PASP automatically verifies statistical assumptions:
\begin{itemize}
    \item \textbf{Normality:} Shapiro-Wilk test.
    \item \textbf{Homogeneity:} Levene's test.
\end{itemize}
If assumptions are violated, it intelligently switches to \textbf{Welch's} corrections or \textbf{Non-parametric} alternatives (Mann-Whitney U, Wilcoxon, Kruskal-Wallis).

\section{Command Reference}

\subsection{Descriptive Statistics}
Classifies variables as Nominal, Ordinal, or Scale.
\begin{itemize}
    \item \texttt{--descriptives}: Flag for Style A.
    \item \texttt{descriptives}: Subcommand for Style B.
    \item \texttt{--vars}: Specify variables to analyze.
\end{itemize}
\textbf{Example:}
\begin{lstlisting}[language=bash]
pasp examples/ejemplo_anova.csv --descriptives --vars score age
\end{lstlisting}

\subsection{One-sample T-Test}
Compares a sample mean against a test value. Uses \textbf{Wilcoxon} if normality fails.
\begin{itemize}
    \item \texttt{--ttest-one VAR VALUE}: Flag for Style A.
    \item \texttt{ttest-one file VAR VALUE}: Subcommand for Style B.
\end{itemize}
\textbf{Example:}
\begin{lstlisting}[language=bash]
pasp examples/ejemplo_ttest_one.csv --ttest-one iq_score 100
\end{lstlisting}

\subsection{Independent Samples T-Test}
Compares means between two independent groups. Fallbacks to \textbf{Welch} or \textbf{Mann-Whitney U}.
\begin{itemize}
    \item \texttt{--ttest-ind VAR GROUP}: Flag for Style A.
    \item \texttt{ttest-ind file VAR GROUP}: Subcommand for Style B.
    \item \texttt{--plot}: Renders a comparative distribution histogram.
\end{itemize}
\textbf{Example:}
\begin{lstlisting}[language=bash]
pasp examples/ejemplo_ttest.csv --ttest-ind recovery_time treatment --plot
\end{lstlisting}

\subsection{Paired Samples T-Test}
\textbf{Example:}
\begin{lstlisting}[language=bash]
pasp examples/ejemplo_ttest_paired.csv --ttest-paired pre_test post_test
\end{lstlisting}

\subsection{One-Way ANOVA}
Includes automatic Bonferroni Post-Hoc comparisons and \textbf{Dunn's test} for non-parametric cases (\textbf{Kruskal-Wallis}).
\begin{itemize}
    \item \texttt{--anova VAR GROUP}: Flag for Style A.
    \item \texttt{anova file VAR GROUP}: Subcommand for Style B.
\end{itemize}
\textbf{Example:}
\begin{lstlisting}[language=bash]
pasp anova examples/ejemplo_anova.csv score age
\end{lstlisting}

\subsection{Correlation and Regression}
\textbf{Examples:}
\begin{lstlisting}[language=bash]
pasp correlation examples/ejemplo_correlacion.csv var1 var2 var3
pasp regression examples/ejemplo_regresion.csv exam_score hours_studied
\end{lstlisting}

\section{Utility Commands}
\begin{itemize}
    \item \texttt{pasp doctor}: Diagnostic report for the environment.
    \item \texttt{pasp examples}: Lists built-in datasets.
    \item \texttt{pasp install-help}: Displays detailed installation guides.
    \item \texttt{pasp --auto [file]}: Quickly runs a summary and descriptives for all numeric data.
\end{itemize}

\section{License}
PASP is released under the MIT License.

\end{document}
